\documentclass[titlepage]{article}
\usepackage[utf8]{inputenc}
\usepackage{amsmath, amssymb, amsthm}
\usepackage{algorithmicx}
\usepackage{graphicx}
\usepackage{subcaption}
\usepackage{listings}
\usepackage{color}
\usepackage{mdframed}
\usepackage{float}
\usepackage{booktabs}
\usepackage{mathtools}

\newcommand{\diag}{\operatorname{diag}}

\definecolor{dkgreen}{rgb}{0,0.6,0}
\definecolor{gray}{rgb}{0.5,0.5,0.5}
\definecolor{mauve}{rgb}{0.58,0,0.82}

\lstset{frame=single,
	language=matlab,
	aboveskip=3mm,
	belowskip=3mm,
	showstringspaces=false,
	columns=flexible,
	basicstyle={\small\ttfamily},
	numbers=none,
	numberstyle=\tiny\color{gray},
	keywordstyle=\color{blue},
	commentstyle=\color{dkgreen},
	stringstyle=\color{mauve},
	breaklines=true,
	breakatwhitespace=true,
	tabsize=3
}

\setlength{\parindent}{0pt}

\title{AAE 50700: Principles of Dynamics\\ Problem Set 5 Report}
\author{Dylan Ashby}
\date{\today}

\begin{document}

\maketitle

\section*{(5.1) Orbital Dynamics with $J_2$ and Atmospheric Drag}

\subsection*{(a): Given}

Consider a spacecraft of mass $m$ orbiting Earth perturbed by $J_2$ and atmosphere drag. Let the inertial frame be $\mathcal{N}:\{\hat{\boldsymbol{e}}_x,\hat{\boldsymbol{e}}_y,\hat{\boldsymbol{e}}_z\}$ and write its Cartesian position as $\boldsymbol{r}=x\hat{\boldsymbol{e}}_x+y\hat{\boldsymbol{e}}_y+z\hat{\boldsymbol{e}}_z,r=\|\boldsymbol{r}\|_2$. Expressing the potential energy for Earth gravity under $J_2$ in the Cartesian coordinates, we have
\begin{equation}
	V(x,y,z)=-\frac{\mu m}{r}+\frac{\mu mJ_2r_o^2}{2r^5}\left(3z^2-r^2\right),
\end{equation}
where $\mu$ is Earth's gravitational parameter and $r_o=R_{\text{earth}}$ is its equatorial radius.\\

We use the same model as considered in PS1 to model the drag, which assumes an inertially stationary, spherically symmetric atmosphere, given by
\begin{equation}
	\boldsymbol{F}_D=m\boldsymbol{a}_D,\quad\boldsymbol{a}_D=-\frac{1}{2}C_D\frac{A}{m}\rho(r)v^2\hat{\boldsymbol{v}},\quad v\coloneqq\|\boldsymbol{v}\|_2,\quad\hat{\boldsymbol{v}}\coloneqq\frac{\boldsymbol{v}}{v},
	\label{eq:dragforce}
\end{equation}
where $\boldsymbol{v}=\dot{\boldsymbol{r}}$, $C_D$ is the drag coefficient, $A$ is the reference area, and
\begin{equation}
	\rho(r) \approx \rho_0 \exp\left[ -\frac{r - 6978~\mathrm{km}}{H} \right].
\end{equation}

\subsubsection*{(a.1): Find}

From the Lagrangian using the Cartesian coordinates as generalized coordinates, $\boldsymbol{q}=[x,y,z]^T$.\\

\textbf{Analysis/Conclusions:}\\

First recall that the kinetic energy $T$ of an object in motion can be expressed by half the mass times the square of the velocity:
\begin{equation}
	T(\boldsymbol{v})=\frac{m}{2}\boldsymbol{v}^T\boldsymbol{v}=\frac{mv^2}{2}
\end{equation}
\begin{mdframed}[linewidth=1pt, innertopmargin=10pt, innerbottommargin=10pt]
	Then, recall the Lagrangian $L$ is the kinetic energy minus the potential energy $L=T-V$ ($V$ given by the problem definition):
	\begin{equation}
		L(\boldsymbol{r},\boldsymbol{v})=\frac{mv^2}{2}+\frac{\mu m}{r}-\frac{\mu mJ_2r_o^2}{2r^5}\left(3z^2-r^2\right)
	\end{equation}
\end{mdframed}

\subsubsection*{(a.2): Find}

Derive the generalized force for Eq.~(\ref{eq:dragforce}) associated with $\boldsymbol{q}=[x,y,z]^T$.\\

\textbf{Analysis/Conclusions:}\\

It is known that virtual work $\delta w$ is translational force $\boldsymbol{F}$ dotted with the Jacobian of position with respect to the generalized coordinates multiplied by the displacement of generalized coordinates. This also equals generalized force $\boldsymbol{Q}$ dotted with the same displacement:
\begin{equation}
	\delta w=\boldsymbol{F}^T\frac{\partial \boldsymbol{r}}{\partial \boldsymbol{q}}\delta\boldsymbol{q}=\boldsymbol{Q}^T\delta\boldsymbol{q}.
\end{equation}

Since in this case the generalized coordinates are the Cartesian position, $\boldsymbol{q}=\boldsymbol{r}$, the Jacobian is the identity matrix, and the generalized force is simply the applied force:
\begin{equation}
	\delta w = \boldsymbol{F}^T\frac{\partial\boldsymbol{r}}{\partial\boldsymbol{r}}\delta\boldsymbol{r}=\boldsymbol{F}^T I\,\delta\boldsymbol{r}=\boldsymbol{F}^T\delta\boldsymbol{r}\quad\Rightarrow\quad\boldsymbol{Q}=\boldsymbol{F}.
\end{equation}

\begin{mdframed}[linewidth=1pt, innertopmargin=10pt, innerbottommargin=10pt]
	So, from Eq.~(\ref{eq:dragforce}) there is only one force, which is the force of drag $\boldsymbol{F}_D$. This means the generalized force from that equation associated with $\boldsymbol{q}$ is $\boldsymbol{Q}=\boldsymbol{F}_D$. The mass terms cancel and simplify to get generalized force:
	\begin{equation}
		\boldsymbol{Q}=\boldsymbol{F}_D=m\boldsymbol{a}_D=m\left(-\frac{1}{2}C_D\frac{A}{m}\rho(r)v^2\hat{\boldsymbol{v}}\right)=-\frac{1}{2}C_DA\rho(r)v^2\hat{\boldsymbol{v}}
	\end{equation}
\end{mdframed}

\subsubsection*{(a.3): Find}

Derive the Lagrange's equations and wrtie the result as the Cartesian vector EOM
\begin{equation}
	\ddot{\boldsymbol{r}}=\boldsymbol{a}_{2B}+\boldsymbol{a}_{J2}+\boldsymbol{a}_D
	\label{eq:lagrange}
\end{equation}

We see here that this system has three degrees a freedom with no redundancy in coordinates, so $\boldsymbol{q}$ is a minimal coordinate set, and the system has no constraints on movement. Since there is also a non-conservative force (drag), the general form of Lagrange's equations with non-conservative general force may be used:
\begin{equation}
	\frac{d}{dt}\left(\frac{\partial L}{\partial\dot{\boldsymbol{q}}}\right)-\frac{\partial L}{\partial\boldsymbol{q}}=\boldsymbol{Q}_{\text{nc}}
\end{equation}

And since, $\boldsymbol{q}=\boldsymbol{r}$ and $\boldsymbol{F}_D$ is the only non-conservative general force:
\begin{equation}
	\frac{d}{dt}\left(\frac{\partial L}{\partial\boldsymbol{v}}\right)-\frac{\partial L}{\partial\boldsymbol{r}}=\boldsymbol{F}_{D}
	\label{eq:lagrangeorbitalcart}
\end{equation}

Then we compute the derivatives in Eq.~(\ref{eq:lagrangeorbitalcart}) individually.\\

So first to solve $\frac{\partial L}{\partial\boldsymbol{v}}$:
\begin{equation}
	\begin{split}
		\frac{\partial L}{\partial\boldsymbol{v}}&=\frac{\partial}{\partial\boldsymbol{v}}\left(\frac{mv^2}{2}+\frac{\mu m}{r}-\frac{\mu mJ_2r_o^2}{2r^5}\left(3z^2-r^2\right)\right) \\
		&=\frac{m}{2}\frac{\partial v^2}{\partial\boldsymbol{v}}+\frac{\partial}{\partial\boldsymbol{v}}\left(\frac{\mu m}{r}\right)-\frac{\partial}{\partial\boldsymbol{v}}\left(\frac{\mu mJ_2r_o^2}{2r^5}\left(3z^2-r^2\right)\right) \\
		&=\frac{m}{2}2\boldsymbol{v}^T+0-0=m\boldsymbol{v}^T \\
	\end{split}
\end{equation}

Then to solve $\frac{d}{dt}\left(\frac{\partial L}{\partial\boldsymbol{v}}\right)$:
\begin{equation}
	\frac{d}{dt}\left(\frac{\partial L}{\partial\boldsymbol{v}}\right)=m\frac{d\boldsymbol{v}^T}{dt}=m\boldsymbol{a}^T 
\end{equation}

Then solve $-\frac{\partial L}{\partial\boldsymbol{r}}$:
\begin{equation}
	\begin{split}
		-\frac{\partial L}{\partial\boldsymbol{r}}&=-\frac{\partial}{\partial\boldsymbol{r}}\left(\frac{mv^2}{2}+\frac{\mu m}{r}-\frac{\mu mJ_2r_o^2}{2r^5}\left(3z^2-r^2\right)\right) \\
		&=-\frac{\partial}{\partial\boldsymbol{r}}\left(\frac{mv^2}{2}\right)-\mu m\frac{\partial r^{-1}}{\partial\boldsymbol{r}}+\frac{\partial}{\partial\boldsymbol{r}}\left(\frac{\mu mJ_2r_o^2}{2r^5}\left(3z^2-r^2\right)\right) \\
		&=0-\mu m\frac{\partial r^{-1}}{\partial\boldsymbol{r}}+\frac{\partial}{\partial\boldsymbol{r}}\left(\frac{\mu mJ_2r_o^2}{2r^5}\left(3z^2-r^2\right)\right) \\
	\end{split}
\end{equation}

Then for readability, we do the remaining component derivatives individually:
\begin{equation}
	\frac{\partial r^{-1}}{\partial\boldsymbol{r}}=\frac{\partial r^{-1}}{\partial r}\frac{\partial r}{\partial\boldsymbol{r}}=-r^{-2}\hat{\boldsymbol{r}}^T
\end{equation}

And the other component:
\begin{equation}
	\begin{split}
		&=\frac{\partial}{\partial\boldsymbol{r}}\left(\frac{\mu mJ_2r_o^2}{2r^5}\left(3z^2-r^2\right)\right)\\
		&=\frac{1}{2}\mu mJ_2r_o^2\left(\frac{\partial}{\partial\boldsymbol{r}}\left(3r^{-5}z^2\right)-\frac{\partial r^{-3}}{\partial\boldsymbol{r}}\right) \\
		&=\frac{1}{2}\mu mJ_2r_o^2\left(3\left(\frac{\partial r^{-5}}{\partial\boldsymbol{r}}z^2+r^{-5}\frac{\partial z^2}{\partial\boldsymbol{r}}\right)-\frac{\partial r^{-3}}{\partial r}\frac{\partial r}{\partial\boldsymbol{r}}\right) \\
		&=\frac{1}{2}\mu mJ_2r_o^2\left(3\left(\frac{\partial r^{-5}}{\partial r}\frac{\partial r}{\partial\boldsymbol{r}}z^2+r^{-5}\frac{\partial z^2}{\partial z}\frac{\partial z}{\partial\boldsymbol{r}}\right)-\frac{\partial r^{-3}}{\partial r}\frac{\partial r}{\partial\boldsymbol{r}}\right) \\
		&=\frac{1}{2}\mu mJ_2r_o^2\left(3\left(\frac{\partial r^{-5}}{\partial r}\frac{\partial r}{\partial\boldsymbol{r}}z^2+r^{-5}\frac{\partial z^2}{\partial z}\hat{\boldsymbol{e}}_z^T\frac{\partial\boldsymbol{r}}{\partial\boldsymbol{r}}\right)-\frac{\partial r^{-3}}{\partial r}\frac{\partial r}{\partial\boldsymbol{r}}\right) \\
		&=\frac{1}{2}\mu mJ_2r_o^2\left(3\left((-5r^{-6})(\hat{\boldsymbol{r}}^T)z^2+2r^{-5}z\hat{\boldsymbol{e}}_z^TI\right)-(-3r^{-4}\hat{\boldsymbol{r}}^T)\right) \\
		&=\frac{1}{2}\mu mJ_2r_o^2\left(3\left((-5r^{-6})(r^{-1}\boldsymbol{r}^T)z^2+2r^{-5}z\hat{\boldsymbol{e}}_z^TI\right)-(-3r^{-4})(r^{-1}\boldsymbol{r}^T)\right) \\
		&=\frac{3}{2}\mu mJ_2r_o^2\left(-5r^{-7}z^2\boldsymbol{r}^T+2r^{-5}z\hat{\boldsymbol{e}}_z^T+r^{-5}\boldsymbol{r}^T\right) \\
		&=\frac{3}{2}\mu mJ_2r_o^2r^{-5}\left(\left(1-\frac{5z^2}{r^2}\right)\boldsymbol{r}^T+2z\hat{\boldsymbol{e}}_z^T\right) \\
		&=\frac{3\mu mJ_2r_o^2}{2r^5}\left(\left(1-\frac{5z^2}{r^2}\right)\boldsymbol{r}^T+2z\hat{\boldsymbol{e}}_z^T\right)
	\end{split}
\end{equation}

So then that gives:
\begin{equation}
	\begin{split}
		-\frac{\partial L}{\partial\boldsymbol{r}}&=-\mu m\left(-r^{-2}\hat{\boldsymbol{r}}^T\right)+\frac{3\mu mJ_2r_o^2}{2r^5}\left(\left(1-\frac{5z^2}{r^2}\right)\boldsymbol{r}^T+2z\hat{\boldsymbol{e}}_z^T\right)\\
		&=\frac{\mu m}{r^3}\boldsymbol{r}^T+\frac{3\mu mJ_2r_o^2}{2r^5}\left(\left(1-\frac{5z^2}{r^2}\right)\boldsymbol{r}^T+2z\hat{\boldsymbol{e}}_z^T\right)
	\end{split}
\end{equation}

Which then plugging the above terms back into Eq.~(\ref{eq:lagrangeorbitalcart}) gives:
\begin{equation}
	m\boldsymbol{a}^T+\frac{\mu m}{r^3}\boldsymbol{r}^T+\frac{3\mu mJ_2r_o^2}{2r^5}\left(\left(1-\frac{5z^2}{r^2}\right)\boldsymbol{r}^T+2z\hat{\boldsymbol{e}}_z^T\right)=-\frac{1}{2}C_DA\rho(r)v^2\hat{\boldsymbol{v}}^T
\end{equation}

Which then solved for acceleration $\boldsymbol{a}=\ddot{\boldsymbol{r}}$:
\begin{equation}
	\begin{split}
		m\boldsymbol{a}&=-\frac{\mu m}{r^3}\boldsymbol{r}-\frac{3\mu mJ_2r_o^2}{2r^5}\left(\left(1-\frac{5z^2}{r^2}\right)\boldsymbol{r}+2z\hat{\boldsymbol{e}}_z\right)-\frac{1}{2}C_DA\rho(r)v^2\hat{\boldsymbol{v}}\\
		\boldsymbol{a}&=-\frac{\mu m}{r^3 m}\boldsymbol{r}-\frac{3\mu mJ_2r_o^2}{2r^5 m}\left(\left(1-\frac{5z^2}{r^2}\right)\boldsymbol{r}+2z\hat{\boldsymbol{e}}_z\right)-\frac{1}{2m}C_DA\rho(r)v^2\hat{\boldsymbol{v}} \\
		\boldsymbol{a}&=-\frac{\mu}{r^3}\boldsymbol{r}-\frac{3\mu J_2r_o^2}{2r^5}\left(\left(1-\frac{5z^2}{r^2}\right)\boldsymbol{r}+2z\hat{\boldsymbol{e}}_z\right)-\frac{1}{2m}C_DA\rho(r)v^2\hat{\boldsymbol{v}} \\
	\end{split}
\end{equation}

\begin{mdframed}[linewidth=1pt, innertopmargin=10pt, innerbottommargin=10pt]
	Which means the Cartesian vector EOM are:
	\begin{equation}
		\ddot{\boldsymbol{r}}=-\frac{\mu}{r^3}\boldsymbol{r}-\frac{3\mu J_2r_o^2}{2r^5}\left(\left(1-\frac{5z^2}{r^2}\right)\boldsymbol{r}+2z\hat{\boldsymbol{e}}_z\right)-\frac{1}{2m}C_DA\rho(r)v^2\hat{\boldsymbol{v}}
	\end{equation}
	Where:
	\begin{equation}
		\begin{split}
			\boldsymbol{a}_{2B}&=-\frac{\mu}{r^3}\boldsymbol{r} \\
			\boldsymbol{a}_{J2}&=-\frac{3\mu J_2r_o^2}{2r^5}\left(\left(1-\frac{5z^2}{r^2}\right)\boldsymbol{r}+2z\hat{\boldsymbol{e}}_z\right) \\
			\boldsymbol{a}_D&=-\frac{1}{2m}C_DA\rho(r)v^2\hat{\boldsymbol{v}} \\
		\end{split}
	\end{equation}
	Since these terms come from the two-body problem, $J_2$ perturbation, and drag acceleration components of the system respectively.
\end{mdframed}




\end{document}